\section{Related Work}

\textbf{Entity matching foundations.} Entity matching has a long history in statistics and data integration, from probabilistic record linkage \cite{fellegi1969record} and duplicate detection surveys \cite{elmagarmid2007duplicate} to book-length treatments of data matching and entity resolution \cite{christen2012data,christophides2015webdata}. Recent surveys frame EM as an end-to-end process for large-scale data integration rather than a standalone binary classifier \cite{christophides2020overview}. Magellan follows this view by supporting a full EM workflow with blocking, labeling, debugging, and evaluation \cite{konda2016magellan}.

\textbf{Neural and PLM-based matchers.} Deep learning broadened the design space for EM by learning representations of serialized records and record pairs \cite{mudgal2018deep,barlaug2021survey}. Transformer-based matchers and PLM-based systems then became strong baselines for benchmark EM tasks \cite{brunner2020transformer,li2020ditto}. Ditto is especially relevant for this paper because it is a widely used downstream matcher for training-data experiments and supports data augmentation and knowledge injection \cite{li2020ditto}. The limitation shared by these supervised systems is their dependence on task-specific labels, which motivates both LLM-based direct matching and LLM-assisted construction of labeled training data.

\textbf{LLMs for entity matching.} Recent work studies whether LLMs can replace or complement PLM-based EM systems. Peeters, Steiner, and Bizer evaluate hosted and open-source LLMs for zero-shot, few-shot, rule-augmented, and fine-tuned EM, showing high zero-shot performance, strong robustness to unseen entities, but also model- and dataset-specific prompt sensitivity \cite{peeters2025llmem}. Their work also demonstrates structured explanations and LLM-assisted error analysis. Steiner, Peeters, and Bizer further investigate fine-tuning LLMs for EM, showing that fine-tuning can improve smaller models and in-domain transfer, while cross-domain transfer and example selection effects are mixed \cite{steiner2025finetuning}. The present paper builds on those findings but changes the unit of analysis: instead of asking only whether LLMs can classify pairs or be fine-tuned for EM, it asks whether selected and LLM-labeled pairs form a reliable, auditable training dataset.

\textbf{Weak supervision and data programming.} Auto-labeling is related to weak supervision, where noisy labeling functions are combined to create training data \cite{ratner2017snorkel}. In EM, weak supervision can reduce manual labeling effort, but it also makes label provenance and noise modeling central concerns. The difference in this paper is that the labeling source is an LLM-driven EM pipeline and the focus is on auditing selected pair labels against benchmark structure, including split overlap, novelty, coverage, and consistency proxies.

\subsection{Active learning for EM}
